\subsection{Độ nhạy cảm với tỉ lệ tập huấn luyện và tập kiểm tra}

\subsubsection{Hồi quy}

\figurename~\ref{fig:regression-split-sens-mse}
và~\ref{fig:regression-split-sens-mae} trình biểu đồ hộp các kết quả của 10 lần
chạy ngẫu nhiên với từng tỉ lệ tập huấn luyện/kiểm tra khác nhau trên các chỉ
số tương ứng MSE, MAE.

\begin{figure}[htbp]
  \centering
  \includegraphics[width=0.9\textwidth]{regression-split-sens-mse}
  \caption{Biểu đồ hộp các chỉ số MSE của các mô hình hồi quy đối với các tỉ lệ
  tập huấn luyện và kiểm tra khác nhau}
  \label{fig:regression-split-sens-mse}
\end{figure}

\begin{figure}[htbp]
  \centering
  \includegraphics[width=0.9\textwidth]{regression-split-sens-mae}
  \caption{Biểu đồ hộp các chỉ số MAE của các mô hình hồi quy đối với các tỉ lệ
  tập huấn luyện và kiểm tra khác nhau}
  \label{fig:regression-split-sens-mae}
\end{figure}

Nhìn vào biểu đồ, ta thấy các mô hình hồi quy tuyến tính cơ bản và chính quy
hóa có ổn định khá cao. Tuy nhiên, các mô hình vẫn còn bị ảnh hưởng bởi cách
chia dữ liệu, như đối với chỉ số MSE, có một điểm xuất hiện với giá trị trên 5.

Trong khi đó, mô hình rất ổn định với cả cách chia này, khi mà ở chỉ số MSE,
các điểm chỉ phân tán trong khoảng 0.66--0.8. Với các cách phân chia dữ liệu
cực đoan, kết quả cũng không cho thấy sự xuất hiện các ngoại lai quá lớn. Nhưng
nếu xét trên chỉ số MAE, thì RBF lại không định bằng các mô hình trên khi dải
phân bố rộng hơn đáng kể.

Các mô hình cơ sở phi tuyến Polynomial và Fourier có sự không ổn định khi mà
xuất hiện các điểm ngoại lai rất lớn như (700,000+ ở MSE của Fourier, 3,000+ ở
MSE của Polynomial, các chỉ số MAE cũng rất lớn nếu so sánh các mô hình khác).
Điều này cho thấy sự mất ổn định cũng như dấu hiệu quá khớp ở các mô hình này.

Nhìn chung, RBF là mô hình ổn định nhất khi khó bị ảnh hưởng bởi cách chia dữ
liệu và kết quả ở 2 chỉ số đều rất tốt. Các mô hình hồi quy tuyến tính cơ bản
và chính quy hóa cũng có độ ổn định cao nhưng có phụ thuộc vào chia dữ liệu
nhẹ.

\subsubsection{Phân lớp}

\figurename~\ref{fig:classification-split-sens-acc}
và~\ref{fig:classification-split-sens-f1} trình biểu đồ hộp các kết quả của 10 lần
chạy ngẫu nhiên với từng tỉ lệ tập huấn luyện/kiểm tra khác nhau trên các chỉ
số tương ứng Accuracy và F1 macro.

\begin{figure}[htbp]
  \centering
  \includegraphics[width=0.9\textwidth]{classification-split-sens-acc}
  \caption{Biểu đồ hộp các chỉ số Accuracy của các mô hình phân lớp đối với các
  tỉ lệ tập huấn luyện và kiểm tra khác nhau}
  \label{fig:classification-split-sens-acc}
\end{figure}

\begin{figure}[htbp]
  \centering
  \includegraphics[width=0.9\textwidth]{classification-split-sens-f1}
  \caption{Biểu đồ hộp các chỉ số F1 của các mô hình phân lớp đối với các tỉ lệ
  tập huấn luyện và kiểm tra khác nhau}
  \label{fig:classification-split-sens-f1}
\end{figure}

Đối với chỉ số Accuracy, mô hình Softmax đạt giá trị cao nhất và ổn định nhất,
dao động quanh mức khoảng 0.711--0.714 ở cả các tỉ lệ chia 60/40, 70/30 và
80/20. Mô hình LDA đứng thứ hai với Accuracy thấp hơn một chút (xấp xỉ
0.678--0.681) nhưng cũng rất ổn định giữa các seed. Trong khi đó, Weighted
Softmax có Accuracy thấp hơn đáng kể (khoảng 0.589--0.596), và QDA cho kết quả
thấp nhất (khoảng 0.24) với phương sai nhỏ nhưng hiệu năng kém rõ rệt so với
các mô hình còn lại.

Đối với chỉ số F1-macro, xu hướng có sự khác biệt rõ hơn. Mô hình LDA thể hiện
hiệu năng tốt nhất và ổn định nhất với F1 khoảng 0.506--0.512, phản ánh khả năng
cân bằng tốt giữa precision và recall trên các lớp. Weighted Softmax đứng thứ
hai với F1 khoảng 0.488--0.495; mặc dù Accuracy thấp hơn, mô hình này có xu
hướng cải thiện recall nên F1 tương đối cạnh tranh. Softmax dù có Accuracy cao
nhất nhưng F1 chỉ quanh mức 0.43--0.44, cho thấy mô hình thiên lệch hơn về lớp
chiếm ưu thế. QDA tiếp tục cho kết quả thấp nhất với F1 khoảng 0.25 và không
cải thiện đáng kể theo thay đổi seed hay tỉ lệ chia dữ liệu.

Nhìn chung, có thể rút ra rằng Softmax tối ưu về Accuracy nhưng không tối ưu về
F1-macro, trong khi LDA là mô hình cân bằng tốt nhất giữa hai tiêu chí và ổn
định nhất trên nhiều lần chạy. Weighted Softmax có cải thiện về khả năng thu
hồi (recall) nhưng chưa vượt trội về tổng thể, còn QDA không phù hợp với bài
toán do hiệu năng thấp và thiếu tính phân biệt giữa các lớp.
